\section{Proofs for Section~\ref{sec:qp}}\label{app:qp}

This appendix proves the results of Section~\ref{sec:qp}. The expected-work
theorems for quadratic programs, mixed-integer quadratic programs, and
separable recourse share one argument, given in Appendix~\ref{app:qp:main};
the separable and anisotropic results are derived from it in
Appendix~\ref{app:qp:sep} by listing the inputs that change.

\subsection{Rational normalization}\label{app:qp:normalize}

\begin{proof}[Proof of Lemma~\ref{lem:qp:normalize}]
\emph{Preliminaries.} Test $A\succeq0$ exactly by symmetric elimination; if it
holds, stop. Otherwise $A\ne0$. Let $H=\max_i\sum_j|A_{ij}|\ge\norm A$, let
$d=\rank A\ge1$, and let $q_A$ be the product of the denominators of the
entries of $A$. The coefficient of $t^{n-d}$ in $\det(tI-q_AA)$ is a nonzero
integer equal, up to sign, to $q_A^d$ times the product of the $d$ nonzero
eigenvalues of $A$. Since the other nonzero eigenvalues have absolute value at
most $H$, every nonzero eigenvalue has absolute value at least
\[
 \mu=q_A^{-d}H^{-(d-1)} .
\]
In particular $\mu\le\nu$. Start from $\beta=H$ and halve $\beta$ while
$A+(\beta/2)I\succeq0$. Each successful test is equivalent to $\beta/2\ge\nu$,
so the final value satisfies $\nu\le\beta<2\nu$, after at most
$1+\log_2(H/\mu)$ tests. Compute the orthogonal projector $R$ onto the range of
$A$ as $R=C(C^{\mathsf T} C)^{-1}C^{\mathsf T}$ for a maximal set $C$ of linearly
independent columns of $A$. All these numbers have polynomial bit length.

\emph{Rational Jacobi rotations.} Put $\epsilon=\mu^2/(16n\beta)$ and
$\tau=\epsilon/n$. If $n=1$, no rotation is needed. Starting from $Q=I$
and $\bar A=A$, repeat while some
off-diagonal entry of $\bar A$ exceeds $\tau$ in absolute value: choose a
largest one, $\bar A_{pq}=b'$, with $a'=\bar A_{pp}$ and $c'=\bar A_{qq}$; for a
rational $t$ let $G(t)$ be the rotation in the $(p,q)$ plane with entries
$c_t=(1-t^2)/(1+t^2)$ and $s_t=2t/(1+t^2)$; replace $Q$ by $QG(t)$ and $\bar A$
by $G(t)^{\mathsf T}\bar AG(t)$. The new $(p,q)$ entry is
$\varphi(t)/(1+t^2)^2$ with
$\varphi(t)=b'(1-6t^2+t^4)+2(c'-a')t(1-t^2)$. Here $\varphi(0)=b'$, and at
$t_{\rm end}=-\sgn(b'(c'-a'))/2$ (or $1/2$ if $c'=a'$),
$\varphi(t_{\rm end})=-\frac7{16}b'-\frac34|c'-a'|\sgn(b')$ has the opposite
sign. Bisection on the sign of $\varphi$ finds a dyadic $t$ within
$\tau/(8H)$ of a root. The rotated block is a principal block of an orthogonal
conjugate of $A$, so its norm is at most $H$; the angle $2\arctan t$ has
derivative at most $2$ in $t$, so the new entry has derivative at most $4H$ in
$t$, and its absolute value is at most $\tau/2$.

A rotation in the $(p,q)$ plane preserves $\bar A_{pr}^2+\bar A_{qr}^2$ for
every $r\notin\{p,q\}$, so the off-diagonal energy
$\operatorname{Off}=\sum_{p\ne q}\bar A_{pq}^2$ changes by
$2(b''^2-b'^2)\le-\frac32b'^2$, where $|b''|\le\tau/2<|b'|/2$ is the new entry.
Since $b'^2\ge\operatorname{Off}/(n(n-1))$, each rotation multiplies
$\operatorname{Off}$ by at most $1-3/(2n(n-1))$. Initially
$\operatorname{Off}\le\norm A_F^2\le nH^2$, and the loop can continue only while
$\operatorname{Off}>\tau^2$. Hence there are $O(n^2\log(nH/\tau))$ rotations.
A rotation with $t=m'/2^{\theta}$ has common denominator $2^{2\theta}+m'^2$; the
common denominators of the products multiply, so their bit lengths add, and
exact orthogonality bounds all entries of $Q$ by $1$. All numbers therefore
have polynomial bit length, and $\log(H/\tau)$ is polynomial because $\log\mu$
is. At termination $Q^{\mathsf T} AQ=\diag(a_1,\dots,a_n)+E$ with zero diagonal in
$E$ and $\norm E\le\norm E_F\le n\tau=\epsilon$.

\emph{Selecting the negative directions.} The eigenvalues of $A$ are at most
$-\mu$ ($k$ of them) or nonnegative. By Weyl's inequality and
$\epsilon\le\mu/16$, exactly $k$ of the $a_j$ are below $-\mu/2$. Let $B$ be
the $n\times k$ matrix of the corresponding columns $q_j$ of $Q$, so
$B^{\mathsf T} B=I_k$, and let $\Pi$ be the orthogonal projector onto the span of the
eigenvectors of $A$ with negative eigenvalues. From $Aq_j=a_jq_j+QEe_j$ we get
$\norm{(A-a_jI)q_j}\le\epsilon$. On the range of $I-\Pi$, the matrix
$A-a_jI$ has eigenvalues at least $\mu/2$, so
$\norm{(I-\Pi)q_j}\le2\epsilon/\mu$. For projectors of equal rank,
$\norm{BB^{\mathsf T}-\Pi}_F^2=2\norm{(I-\Pi)B}_F^2$, hence
$\norm{BB^{\mathsf T}-\Pi}\le2\sqrt{2k}\,\epsilon/\mu\le3n\epsilon/\mu$.

\emph{The correction.} Put $U=RB$ and $T=U^{\mathsf T}$. Since the range of $\Pi$ lies
in the range of $R$, $R\Pi R=\Pi$, and
$\norm{UU^{\mathsf T}-\Pi}=\norm{R(BB^{\mathsf T}-\Pi)R}\le3n\epsilon/\mu$. On the range of
$A$, the matrix $A+2\beta\Pi$ has eigenvalues at least $\mu$: a negative
eigenvalue $\lambda\ge-\nu$ becomes $\lambda+2\beta\ge\beta\ge\mu$, and a
positive one is at least $\mu$. Replacing $\Pi$ by $UU^{\mathsf T}$ changes the
matrix by at most $2\beta\cdot3n\epsilon/\mu=3\mu/8$, so
$A+2\beta UU^{\mathsf T}\succeq\frac58\mu I$ on the range of $A$. On $\ker A$, both
$A$ and $U^{\mathsf T}=B^{\mathsf T} R$ vanish. Hence $A+2\beta T^{\mathsf T} T\succeq0$, and the
rows of $T$ lie in the range of $A$.

\emph{Frame.} $\norm U\le\norm R\norm B\le1$, so $TT^{\mathsf T}\preceq I$. Moreover
$TT^{\mathsf T}=B^{\mathsf T} RB=I-B^{\mathsf T}(I-R)B$, and
$\norm{(I-R)B}_F^2\le\norm{(I-\Pi)B}_F^2\le4k\epsilon^2/\mu^2\le k/(64n^2)\le1/64$,
using $\mu\le\beta$. Thus $TT^{\mathsf T}\succeq\frac{63}{64}I$.
\end{proof}

\subsection{The auxiliary value function and face candidates}

\begin{proof}[Proof of Lemma~\ref{lem:qp:aux}]
(a) For every $x$, expanding the square gives
\[
 F_\gamma(x)+\tfrac12(a-Tx+\Lambda^{-1}\xi)^{\mathsf T}\Lambda(a-Tx+\Lambda^{-1}\xi)
 -\tfrac12\xi^{\mathsf T}\Lambda^{-1}\xi
 =F(x)+\zeta^{\mathsf T} x+\tfrac12(a-Tx)^{\mathsf T}\Lambda(a-Tx)+\xi^{\mathsf T} a ,
\]
because $\gamma^{\mathsf T} x=\xi^{\mathsf T} Tx+\zeta^{\mathsf T} x$. Minimizing over $x$ gives
the identity. The right side of the identity is at least
$f^*-\frac12\xi^{\mathsf T}\Lambda^{-1}\xi$, with equality at $a=Tx^*-\Lambda^{-1}\xi$
and $x=x^*$. For an attaining witness $x_a$, dropping the nonnegative square
gives $V(a)\ge F_\gamma(x_a)-\frac12\xi^{\mathsf T}\Lambda^{-1}\xi$.

(b) $W_\zeta(a)-\frac12a^{\mathsf T}\Lambda a=\min_x\bigl[F(x)+\zeta^{\mathsf T} x+\frac12(Tx)^{\mathsf T}\Lambda Tx-a^{\mathsf T}\Lambda Tx\bigr]$
is the infimum of a compactly indexed family of affine functions of $a$,
and is therefore concave. Compactness and continuity separately give
attainment; the family need not be finite.

(c) For every $z$, the witness $x_a$ is feasible in the problem defining
$W_\zeta(z)$, so
$W_\zeta(z)\le W_\zeta(a)+(a-Tx_a)^{\mathsf T}\Lambda(z-a)+\frac12(z-a)^{\mathsf T}\Lambda(z-a)$.
Adding $\xi^{\mathsf T} z$ gives the upper model. Evaluating it at
$z=a-\Lambda^{-1}g$ and using $V(z)\ge V^*$ gives the inequality.

(d) The minimizer $Tx^*-\Lambda^{-1}\xi$ from (a) lies in $\mathcal A$.
\end{proof}

\begin{proof}[Proof of Lemma~\ref{lem:qp:face}]
(a) A global minimizer exists by compactness. Among all global minimizers
choose $\hat x$ whose minimal face $\Phi$ (the face containing $\hat x$ in its
relative interior) has least dimension. Let $\mathcal S$ be the set of rows
active at $\hat x$; the affine hull of $\Phi$ is $\{x:D_{\mathcal S}x=e_{\mathcal S}\}$,
with tangent space $\ker D_{\mathcal S}$. For $v\in\ker D_{\mathcal S}$ and small
$|t|$, $\hat x+tv\in\mathcal Q$, so $\nabla G(\hat x)^{\mathsf T} v=0$ and
$v^{\mathsf T} Hv\ge0$. If $v\ne0$ and $v^{\mathsf T} Hv=0$, then $G$ is constant on the line
$\hat x+tv$, which leaves the bounded face $\Phi$ through its relative
boundary at a global minimizer on a face of smaller dimension, a
contradiction. Hence $H$ is positive definite on $\ker D_{\mathcal S}$, or
$\ker D_{\mathcal S}=\{0\}$. Let $S\subseteq\mathcal S$ index a maximal linearly
independent set of rows, so $\ker D_S=\ker D_{\mathcal S}$. If
$K_S(v;\mu)=0$, then $D_Sv=0$ and $v^{\mathsf T} Hv=-v^{\mathsf T} D_S^{\mathsf T}\mu=0$, so $v=0$,
and then $D_S^{\mathsf T}\mu=0$ gives $\mu=0$. Finally $\nabla G(\hat x)$ is orthogonal
to $\ker D_S$, hence equals $-D_S^{\mathsf T}\mu$ for some $\mu$, and
$D_S\hat x=e_S$.

(b) Every retained candidate is feasible, so the least retained value is at
least the minimum, and by (a) the minimizer $\hat x$ is retained. Each subset
costs polynomial rational work, and Cramer's rule bounds the bit lengths of
the solutions.

(c) Each nonempty slice is a bounded polytope on which the objective is a
quadratic in the continuous variables; apply (b) to every slice and compare.

(d) The faces of a box are obtained by fixing each coordinate at its lower
bound, fixing it at its upper bound, or leaving it free, and the rows active
on a face form an independent set unless a coordinate has equal bounds.
\end{proof}

\begin{proof}[Proof of Lemma~\ref{lem:qp:growth-sections}]
Fix $\varepsilon$, the coordinate $i$, and the other coefficients, and write
$F_t$ for the objective with $\gamma_i=t$. Call a \emph{candidate type} a pair
$\omega=(z,S)$ of a label and a set of linearly independent rows of its slice.
For a fixed Hessian, whether the stationarity matrix of $\omega$ is
nonsingular does not depend on $t$; there are at most $\mathsf F$ types. For
every type with a nonsingular matrix for the Hessian $A$, the candidate
$x_\omega(t)$ is affine in $t$, and the set of $t$ where it is feasible is an
interval. For such $\omega$, the modified quadratic
$\tilde F_{\omega,t}(y)=F_t(y)-\varepsilon\norm{y-x_\omega(t)}^2$ has Hessian
$A-2\varepsilon I$ and a linear coefficient affine in $t$. Its candidates
$y_{\omega,\omega'}(t)$, for types $\omega'$ with nonsingular matrices for
$A-2\varepsilon I$, are affine in $t$ with interval feasibility sets, and
$d_{\omega,\omega'}(t)=\tilde F_{\omega,t}(y_{\omega,\omega'}(t))-F_t(x_\omega(t))$
is a polynomial of degree at most two.

We claim that $g_*\ge\varepsilon$ if and only if some $x_\omega(t)$ is
feasible and $d_{\omega,\omega'}(t)\ge0$ for every $\omega'$ with
$y_{\omega,\omega'}(t)$ feasible. If this holds, then by
Lemma~\ref{lem:qp:face}(c) applied to $\tilde F_{\omega,t}$ the minimum of
$\tilde F_{\omega,t}$ over $X$ is attained at a feasible modified candidate, so
$\min_X\tilde F_{\omega,t}\ge F_t(x_\omega(t))=\tilde F_{\omega,t}(x_\omega(t))$,
which is the growth inequality at $x_\omega(t)$. Conversely, if
$g_*\ge\varepsilon$, the minimizer $x^*$ is unique, so by
Lemma~\ref{lem:qp:face} it is a candidate $x_\omega(t)$, and the growth
inequality makes $x^*$ a minimizer of $\tilde F_{\omega,t}$, so all $d\ge0$.

The truth value of this condition can change only at the at most $2\mathsf F$
endpoints of feasibility intervals of the $x_\omega$, the at most
$2\mathsf F^2$ endpoints for the $y_{\omega,\omega'}$, and the at most
$2\mathsf F^2$ roots of nonzero polynomials $d_{\omega,\omega'}$. With at most
$K_{\rm b}=2\mathsf F+4\mathsf F^2$ breakpoints, the good set and the bad set
are unions of at most $2K_{\rm b}+1\le8(\mathsf F+1)^2$ intervals. The
probability bounds follow from Theorem~\ref{thm:count:growth-tail} and
Lemma~\ref{lem:count:transfer}(b), with Kolmogorov distances $1/M$ and
$2^{-b}$ and density bounds $1/(2\sigma)$ and $1/(\sigma\sqrt{2\pi})$.
\end{proof}

\subsection{The growth-conditioned algorithm and two negative directions}

\begin{proof}[Proof of Theorem~\ref{thm:qp:conditioned}]
\emph{Algorithm.} Compute $\beta$ and $T$ by Lemma~\ref{lem:qp:normalize} and
put $\alpha=2\beta$, $\Lambda=\alpha I$, $\zeta=0$, $\xi=0$, so that
$V=W_0$ and $V^*=f^*$. Compute the ranges $[\ell_i,u_i]$ of $(Tx)_i$ over $X$
by linear programming, and remove the coordinates with $\ell_i=u_i$ only
from the mesh. Every reduced auxiliary query is embedded in the full
auxiliary space with these coordinates fixed at their original values;
all rows of $T$ and the full certified PSD matrix $P$ remain in the inner
problem. If no mesh coordinate remains, one exact convex solve at the
constant value of $Tx$ returns a
minimizer by Lemma~\ref{lem:qp:aux}(a). By Lemma~\ref{lem:qp:face}(a) and
Cramer's rule, there are integers $\mathsf V$ and $\mathsf R$, with
polynomially many bits and computable from the data, such that $f^*$ has a
denominator at most $\mathsf V$ and, if the minimizer $x^*$ is unique, every
coordinate of $Tx^*$ has a denominator at most $\mathsf R$. Run the search of
Definition~\ref{def:count:search} on $W_0$ over the box of ranges with
isotropic meshes and curvature $\alpha$, storing an exact attaining witness at
every corner. At the end of every level $j$ do the following.
\begin{enumerate}[label=(\arabic*)]
\item If no cell is retained, return the witness of the incumbent.
\item If $E_j<1/(2\mathsf V^2)$, find the unique rational $\hat f$ with
 denominator at most $\mathsf V$ in the interval $[\underline U_j,U_j]$ of
 Theorem~\ref{thm:count:cells}(v) by continued fractions
 \cite{grotschel1988-geometric-algorithms-and-combinatorial-optimization}. If a
 stored witness $x$ has $F(x)=\hat f$, return $x$.
\item Only after step (2) has reconstructed $\hat f$, with $v_j$ a corner
 attaining $U_j$, look in each remaining mesh coordinate for a
 rational with denominator at most $\mathsf R$ within $1/(4\mathsf R^2)$ of
 $(v_j)_i$. If all exist, solve the inner problem exactly at the resulting
 point $\hat a$, embedded with all removed coordinates fixed, and return
 its witness $\hat x$ if $F(\hat x)=\hat f$.
\end{enumerate}

\emph{(a) Correctness.} In (1), Theorem~\ref{thm:count:cells}(ii) and
Lemma~\ref{lem:qp:aux}(a) show that the witness is a minimizer. In (2) and (3),
the interval $[\underline U_j,U_j]$ contains $f^*$ and has length less than
$1/\mathsf V^2$, while two distinct rationals with denominators at most
$\mathsf V$ differ by at least $1/\mathsf V^2$; hence $\hat f=f^*$, and any
returned feasible point of value $\hat f$ is a minimizer. This holds on every
instance.

\emph{(b) Termination.} Let $g>0$ and $t^*=Tx^*$. For every $a$ and feasible
$x$, $\norm{a-t^*}\le\norm{a-Tx}+\norm{x-x^*}$ since $\norm T\le1$, and the
Cauchy--Schwarz inequality gives
$\frac\alpha2\norm{a-Tx}^2+g\norm{x-x^*}^2\ge g_W\norm{a-t^*}^2$ with
$g_W=\alpha g/(2g+\alpha)$. Hence $W_0(a)-f^*\ge g_W\norm{a-t^*}^2$ and
$\alpha/g_W=2+\alpha/g<\kappa$. By Theorem~\ref{thm:count:cells}(iv), every
cell retained at level $j$ has a corner $v$ with $W_0(v)\le f^*+2E_j$, so
$\norm{v-t^*}^2\le2E_j/g_W\le k\kappa h_j^2/4$. In a refined coordinate the
grid spacing exceeds $h_j/2$, so the box of half-width $h_j\sqrt{k\kappa}/2$
around $t^*$ contains at most $2\sqrt{k\kappa}+2$ grid values per coordinate.
Hence at most $2^k(2\sqrt{k\kappa}+2)^k$ cells are retained and at most
$4^k(2\sqrt{k\kappa}+2)^k$ are processed at the next level. Step (2) applies
once $E_j\le k\alpha s^24^{-j}/8<1/(2\mathsf V^2)$, after $\poly(L_{\rm in})$
levels. If no stored witness attains $f^*$, then
$W_0(v_j)=U_j\le f^*+E_j$ by Theorem~\ref{thm:count:cells}(iii), so
$\norm{v_j-t^*}^2\le E_j/g_W\le k\kappa s^24^{-j}/8$. Once this is less than
$1/(16\mathsf R^4)$, which happens after $\poly(L_{\rm in})+\frac12\log_2\kappa$
levels, each coordinate of $t^*$ is the unique rational with denominator at
most $\mathsf R$ within $1/(4\mathsf R^2)$ of $(v_j)_i$, so $\hat a=t^*$. At
$t^*$ every inner minimizer $\hat x$ satisfies
$F(\hat x)+\frac\alpha2\norm{t^*-T\hat x}^2=W_0(t^*)=f^*$, so $F(\hat x)=f^*$ and
step (3) returns. Level-$j$ corners have bit length polynomial in
$L_{\rm in}+j$, so each exact convex solve, reconstruction, and comparison
costs $\poly(L_{\rm in}+j)$. Summing over the levels proves (b).

\emph{(c)} For $k\le2$ and $Z\ge1$, $\kappa\le2+4Z\le6Z$, so
$(1+\sqrt\kappa)^k\le(1+\sqrt6)^2Z^{k/2}$ and
$J_g\le\poly(L_{\rm in})+3+\log_2Z$. Substituting in (b) gives the bound.
\end{proof}

\begin{proof}[Proof of Corollary~\ref{cor:qp:conditioned-mixed}]
Use the same full convexifying factorization and auxiliary objective, but
solve every inner query by Proposition~\ref{prop:qp:primitives}(b) and
polish its returned integer label by part (a). Compute auxiliary coordinate
ranges over the rational relaxation $\mathcal P$, so they enclose $TX$ and
are available by polynomial-time LP. Remove singleton auxiliary ranges
only from the mesh, with the same full-space embedding as above.
The mixed smallest-face argument of Lemma~\ref{lem:qp:face}(c) and the
LP bounds on integer labels give precomputable polynomial-height
denominator bounds for $f^*$ and, under uniqueness, $Tx^*$.
The growth-transfer inequality holds for every $x\in X$ without using
convexity of $X$. Upper coordinate curvature, corner rounding, the
near-optimal-corner packing, and the two guarded rational reconstructions
therefore apply verbatim. At $Tx^*$ every exact mixed recourse witness has
original objective value $f^*$ by square completion. All later queried and
polished data have polynomial length in $L_{\rm in}+j$; only the inner
solve costs the additional factor $f(n_z)$. This proves the stated work
and correctness contracts.
\end{proof}

\begin{proof}[Proof of Theorem~\ref{thm:qp:two}]
Every sampled coefficient has bit length $O(\log M)$ in case (a) and
$O(b)$ in case (b), plus the bit length of $\sigma$. The least sufficient
choices, or the stated polynomial-bit enlargements, therefore give
$L_{\rm in}\le\poly(I)$, including the sampler's polynomial cost.
If $S=0$, LP fixes every original coordinate and direct evaluation returns
the rational point and value; no growth threshold or mesh is used. Below
assume $S>0$ and $k\ge1$; $k=0$ uses the convex primitive.
Interleave, one bit operation at a time, the algorithm of
Theorem~\ref{thm:qp:conditioned} and the fallback of Lemma~\ref{lem:qp:face}(b)
on the same sampled objective, and stop when either returns. Both are exact,
so the output is exact on every draw. The work is at most twice the smaller
of the two running times, up to a polynomial simulation overhead. The
fallback costs at most $B\poly(L_{\rm in})$. If $g_*>0$, the conditioned
algorithm costs at most $\poly(L_{\rm in})Z^p(1+\log Z)^{d_0}$ with $p=k/2$ and
$Z=\max\{1,\nu/g_*\}$. Put $Z=\infty$ if $g_*=0$. If $Z^p\le B$, then
$\log Z\le p^{-1}\log B$; otherwise the minimum below is at most $B$. Hence
\[
 \min\{B,Z^p(1+\log Z)^{d_0}\}\le(1+p^{-1}\log B)^{d_0}\min\{B,Z^p\}.
\]
For $t\ge1$, $\{Z>t\}=\{g_*<\nu/t\}$. By Lemma~\ref{lem:qp:growth-sections}
with $\mathsf F=2^r$, $\Prob\{Z>t\}\le r_0/t+\beta$ with $r_0=\nu S/\sigma$ and
$\beta=2n\mathsf D/M\le1/B$ in case (a), and with $r_0=\sqrt{2/\pi}\,\nu S/\sigma$
and $\beta=2n\mathsf D2^{-b}\le1/B$ in case (b). Since $\min\{B,Z^p\}\ge1$,
\[
 \E\min\{B,Z^p\}=1+\int_1^B\Prob\{Z>t^{1/p}\}\,dt
 \le1+r_0\int_1^Bt^{-1/p}\,dt+\beta B,
\]
which is at most $2+r_0$ for $p=1/2$ and at most $2+r_0\log B$ for $p=1$.
Since $\log B=O(r+1)$, the expected work is at most
$\poly(I)(1+\nu S/\sigma)$. For $k=0$, $A\succeq0$ and one exact convex solve
suffices.
\end{proof}

\subsection{Pieces, gaps, isolation, and tubes}

\begin{proof}[Proof of Lemma~\ref{lem:qp:pieces}]
Write $x=(z,x_c)$, and let $P_{cc}$, $D_c$, $T_c$ be the blocks of $P$, $D$, $T$
for the continuous variables. For fixed $(a,\zeta)$ the slice objective is
$\frac12x_c^{\mathsf T} P_{cc}x_c+p(a,\zeta)^{\mathsf T} x_c$ plus a constant, where
$p(a,\zeta)=P_{cz}z+b_c+\zeta_c-\alpha T_c^{\mathsf T} a$ is affine, and the slice is
$\mathcal Q_z=\{x_c:D_cx_c\le e-D_zz\}$. Since $P\succeq0$, $P_{cc}\succeq0$.

\emph{Extraction.} Solve the slice exactly
(Proposition~\ref{prop:qp:primitives}) and let $x^0$ be the optimal point and
$g^0=P_{cc}x^0+p(a,\zeta)$. For $y\in\mathcal Q_z$ the objective gap equals
$g^{0\mathsf T}(y-x^0)+\frac12(y-x^0)^{\mathsf T} P_{cc}(y-x^0)$, a sum of two nonnegative
terms, so the optimal set is the polytope
$\mathcal O=\mathcal Q_z\cap\{P_{cc}(y-x^0)=0,\ g^{0\mathsf T}(y-x^0)=0\}$. Find a vertex
$y$ of $\mathcal O$ by linear programming. Let $\mathcal S$ be the rows active at
$y$. If $v\in\ker D_{c,\mathcal S}$ with $v^{\mathsf T} P_{cc}v=0$, then $P_{cc}v=0$,
$y\pm tv\in\mathcal Q_z$ for small $t$, optimality gives $g^{0\mathsf T}v=0$ (the
gradient at $y$ equals $g^0$ because $P_{cc}(y-x^0)=0$), and $y\pm tv\in\mathcal O$,
contradicting that $y$ is a vertex. Hence $P_{cc}$ is positive definite on
$\ker D_{c,\mathcal S}$. By the normal-cone formula for polyhedra there are
multipliers $\mu\ge0$ with $-(P_{cc}y+p)=D_{c,\mathcal S}^{\mathsf T}\mu$; find them by
linear programming. While the rows in the support of $\mu$ are dependent, move
$\mu$ along a dependency until a support coefficient vanishes, keeping
$\mu\ge0$. Extend the final independent support to a maximal independent set
$S\subseteq\mathcal S$, with zero multipliers on the added rows. Then
$\ker D_{c,S}=\ker D_{c,\mathcal S}$, so
$K_S=\bigl(\begin{smallmatrix}P_{cc}&D_{c,S}^{\mathsf T}\\D_{c,S}&0\end{smallmatrix}\bigr)$
is nonsingular, as in the proof of Lemma~\ref{lem:qp:face}.

\emph{(a)--(b).} The system $K_S(x_c;\lambda)=(-p(a',\zeta');e_S-D_{z,S}z)$ has
an affine solution in $(a',\zeta')$; the derivative in $a'$ is the first block
$\Xi_S$ of $K_S^{-1}(\alpha T_c^{\mathsf T};0)$, which depends on neither $\zeta'$, $z$,
$b$, nor $e$. At the query, $(y,\mu_S)$ solves the system, so $a$ lies in the
region for the queried $\zeta$. The same vertex and multiplier existence
argument applies to every real $(a,\zeta)$ for each feasible fixed label;
it requires no rational computation for this existence assertion. Thus the
finitely many regions cover that label's full real parameter space,
including the conditioned lines used in the section-count proof.
On the region the point is feasible and satisfies the
Karush--Kuhn--Tucker conditions with nonnegative multipliers on the active
rows $S$; for a convex quadratic program these conditions are sufficient. The
region is described by the $r$ feasibility rows and the $|S|\le n$ sign
conditions, with normals formed from $\Xi_S$ and the $a'$-derivative of
$\lambda$.

\emph{(c).} Let $\Phi(a',x)=F(x)+\zeta'^{\mathsf T} x+\frac\alpha2\norm{a'-Tx}^2$. Then
$\partial_{a'}\Phi=\alpha(a'-Tx)$ and $\nabla_{x_c}\Phi=P_{cc}x_c+p(a',\zeta')$,
which equals $-D_{c,S}^{\mathsf T}\lambda$ along the solution. Since
$D_{c,S}x_{z,S}(a',\zeta')=e_S-D_{z,S}z$ for all $a'$, $D_{c,S}\Xi_S=0$, and the
chain rule gives
$\nabla_{a'}q_{z,S}=\alpha(a'-Tx_{z,S})-(D_{c,S}\Xi_S)^{\mathsf T}\lambda=\alpha(a'-Tx_{z,S})$.
Hence $H_S=\alpha(I-T_c\Xi_S)$. Multiply $K_S$ and the right side
$(\alpha T_c^{\mathsf T};0)$ by a common denominator of $P$, $D$, and $\alpha T^{\mathsf T}$.
By Cramer's rule every entry of $\Xi_S$ is a quotient of two integer
determinants of order at most $2n$, the denominator nonzero, with entries
bounded by $C_0$; hence its absolute value is at most $U_0$. Every entry of
$T_c\Xi_S$ is then at most $n\tau_TU_0$, and
$\norm{T_c\Xi_S}\le\norm{T_c\Xi_S}_F\le kn\tau_TU_0$. There are at most $Z$
labels and $2^r$ row sets.
\end{proof}

\begin{proof}[Proof of Lemma~\ref{lem:qp:gap}]
(a) Distinct integer vectors differ in some coordinate by at least one, so the
union of the $2n_z$ restricted problems consists exactly of the points of $X$
with a label other than $z_*$. Each restricted problem is a convex MIQP with
the same integer variables and one more row.

(b) For a label $z$ put $m_z(a)=W^z_\zeta(a)-\frac\alpha2\norm a^2$, the
minimum over the points $x$ of label $z$ of affine functions of $a$ with
slopes $-\alpha Tx$, $Tx\in T\mathcal P$. For $a$ and $v$, using the minimizer
at $v$ as a test point and the bound on all slopes,
\[
 -\alpha\max_{y\in T\mathcal P}(a-v)^{\mathsf T} y\le m_z(a)-m_z(v)\le-\alpha\min_{y\in T\mathcal P}(a-v)^{\mathsf T} y .
\]
The interval is the same for every label and has length at most
$\alpha\norm{a-v}\diam(T\mathcal P)\le\alpha\sqrt k\,h\,\bar w$ for $a\in C$.
Hence, for $z\ne z_*$ and $a\in C$,
$W^z_\zeta(a)-W^{z_*}_\zeta(a)\ge\Delta(v)-\alpha\sqrt kh\bar w\ge0$. The label
$z_*$ is therefore optimal on $C$, and on $C\subseteq\mathcal R_{z_*,S}(\zeta)$
its value is $q_{z_*,S}$ with attaining point $x_{z_*,S}$ by
Lemma~\ref{lem:qp:pieces}.

(c) For every feasible $x$ and every $a$, Lemma~\ref{lem:qp:aux}(a) gives
$F_\gamma(x)-f^*\le\bigl[F(x)+\zeta^{\mathsf T} x+\frac\alpha2\norm{a-Tx}^2+\xi^{\mathsf T} a\bigr]-V^*$.
For the winning witness at $v$ the bracket is $V(v)$, and for the best
witness of another label it is $V(v)+\Delta(v)$. Both are at most
$V^*+2E+\theta_{\rm gap}h\le V^*+(k\alpha/4+\theta_{\rm gap})h$, using
$2E\le k\alpha h^2/4\le k\alpha h/4$.
\end{proof}

\begin{proof}[Proof of Lemma~\ref{lem:qp:isolation}]
Fix $i$ and all coefficients except $\gamma_i$. For each value $t$ taken by
$z_i$ on $\mathcal Z$, put
$a_t=\min\{h(z)+\sum_{l\ne i}\gamma_lz_l:z\in\mathcal Z,\ z_i=t\}$. As a
function of $u=\gamma_i$, the least value is the lower envelope of the lines
$a_t+tu$, whose slopes are distinct integers in $[L_i,U_i]$. The active slope
is nonincreasing in $u$ and no slope is active on two separated intervals, so
the envelope has at most $U_i-L_i$ breakpoints, which are fixed under the
conditioning. Suppose that at $u$ some line $a_t+tu$ attains the envelope and
another line $a_{t'}+t'u$ lies within $\varepsilon$ of it. The two lines
cross within distance $\varepsilon/|t-t'|\le\varepsilon$ of $u$. Moving from
$u$ towards the crossing, the envelope either changes its active line before
the crossing or has both lines active at the crossing, which is a breakpoint
because their slopes differ. Hence $u$ lies within $\varepsilon$ of a
breakpoint, an event of conditional probability at most
$(U_i-L_i)\pi_i(2\varepsilon)$.

If $\Delta\le\varepsilon$, a best label and a different second-best label
differ in some coordinate $i$, and their groups satisfy the preceding event
for that $i$. A union over $i$, followed by integration over the other
coefficients, proves the bound.
\end{proof}

\begin{proof}[Proof of Lemma~\ref{lem:qp:tube}]
\emph{The family.} For every pair $\omega=(z,S)$ write
$\nabla_aq_\omega(a,\zeta)=H_\omega a+p_\omega(\zeta)$, with $p_\omega$ affine.
If $H_\omega$ is nonsingular, then for every region row $N a\le n_0+\Gamma\zeta$
with $N\ne0$ let $u=H_\omega^{-\mathsf T}N^{\mathsf T}/\norm{H_\omega^{-\mathsf T}N^{\mathsf T}}$ and
\[
 \ell(\xi,\zeta)=\frac{N H_\omega^{-1}(\xi+p_\omega(\zeta))+n_0+\Gamma\zeta}{\norm{H_\omega^{-\mathsf T}N^{\mathsf T}}} .
\]
For fixed $\zeta$, $|\ell(\xi,\zeta)|$ is the distance from $\xi$ to the image
of the hyperplane $\{Na=n_0+\Gamma\zeta\}$ under the affine map
$\Phi_\zeta(a)=-(H_\omega a+p_\omega(\zeta))$. If $H_\omega$ is singular, choose a
unit vector $u\in\ker H_\omega$ (a Hessian is symmetric, so $u^{\mathsf T}H_\omega=0$), and let
$\ell(\xi,\zeta)=u^{\mathsf T}(\xi+p_\omega(\zeta))$. Each pair contributes at most
$r+n+1$ functionals.

\emph{The estimate.} Let $\omega=(z_v,S_v)$ and $g=\nabla_aq_\omega(v,\zeta)+\xi$.
By Lemma~\ref{lem:qp:pieces}(c), $g=\alpha(v-Tx_\omega(v,\zeta))+\xi$, where
$x_\omega(v,\zeta)$ is an attaining witness of $W_\zeta(v)$ because $z_v$ wins.
Lemma~\ref{lem:qp:aux}(c) gives
$\norm g^2\le2\alpha(V(v)-V^*)\le4\alpha E_j\le k\alpha^2h_j^2/2$. If $H_\omega$
is singular, $u^{\mathsf T}\nabla_aq_\omega(v,\zeta)=u^{\mathsf T} p_\omega(\zeta)$, so
$|\ell(\xi,\zeta)|=|u^{\mathsf T} g|\le\norm g$. If $H_\omega$ is nonsingular, some
corner $a'$ of the cell violates a region row that $v$ satisfies. The normal of
that row is nonzero, since a row with zero normal holds at $v$ exactly when it
holds everywhere. The segment from $v$ to $a'$ meets the boundary hyperplane
at a point $a_\Pi$ with $\norm{a_\Pi-v}\le\diam C\le\sqrt kh_j$. Then
\[
 |\ell(\xi,\zeta)|\le\norm{\xi-\Phi_\zeta(a_\Pi)}
 \le\norm{\xi-\Phi_\zeta(v)}+\norm{H_\omega(v-a_\Pi)}
 \le\norm g+\Theta\sqrt kh_j ,
\]
since $\xi-\Phi_\zeta(v)=g$. In both cases
$|\ell|\le\sqrt k(\alpha/\sqrt2+\Theta)h_j$.

\emph{Ambient form.} With $\xi=\Psi\gamma$ and $\zeta=\Pi_T\gamma$,
$\ell=L^{\mathsf T}\gamma+c$ with $L=\Psi^{\mathsf T} u+\Pi_Tw$, where $w$ collects the
coefficients of $\zeta$. Since $T\Psi^{\mathsf T}=TT^{\mathsf T}(TT^{\mathsf T})^{-1}=I$ and
$T\Pi_T=0$, $TL=u$, so $\norm L\ge\norm u/\norm T\ge1$.
\end{proof}

\begin{lemma}[Tube probabilities]\label{lem:qp:tube-prob}
Let $\ell(\xi,\zeta)=u^{\mathsf T}\xi+w^{\mathsf T}\zeta+c$ be one of the functionals of
Lemma~\ref{lem:qp:tube}, and $\eta>0$.
\begin{enumerate}[label=(\alph*)]
\item If $\zeta=0$ and $\xi_1,\dots,\xi_k$ are independent with law
 $U_{\sigma,M}$, then $\Prob\{|\ell|\le\eta\}\le\sqrt k\eta/\sigma+1/M$.
\item If $T$ has full row rank, $\norm T\le1$, and $\gamma_1,\dots,\gamma_n$ are
 independent with law $U_{\sigma,M}$, then
 $\Prob\{|\ell|\le\eta\}\le\sqrt n\eta/\sigma+1/M$.
\item Under the conditions of (b) with independent $\gamma_i$ satisfying
 $d_{\rm K}(\gamma_i,N(0,\sigma^2))\le\delta_{\rm K}$,
 $\Prob\{|\ell|\le\eta\}\le\eta/\sigma+2n\delta_{\rm K}$.
\end{enumerate}
\end{lemma}

\begin{proof}
(a) Some coordinate has $|u_i|\ge k^{-1/2}$. Conditional on the others, the
event confines $\xi_i$ to an interval of length $2\eta/|u_i|\le2\sqrt k\eta$.
(b) In the form $|L^{\mathsf T}\gamma+c|\le\eta$ with $\norm L\ge1$, some coordinate
has $|L_i|\ge n^{-1/2}$, and the same argument applies. (c) Under the proxy
$N(0,\sigma^2I_n)$, $L^{\mathsf T}\gamma$ is normal with standard deviation at least
$\sigma$, so the probability is at most $2\eta/(\sigma\sqrt{2\pi})\le\eta/\sigma$.
Every coordinate section of the event is an interval, possibly empty or the
whole line, so Lemma~\ref{lem:count:transfer}(b) with $C=1$ adds at most
$2n\delta_{\rm K}$.
\end{proof}

\subsection{Section counts and local counts under the three laws}

In the following lemma, a \emph{closure structure} is one of: continuous
pieces ($n_z=0$) or mixed pieces of Lemma~\ref{lem:qp:pieces}, with $R=Z2^r$;
or separable pieces of Section~\ref{sec:qp:separable}, with $R=R_{\rm sep}$.
Put $P_{\rm sec}=2R$ for continuous and separable pieces and
$P_{\rm sec}=2R^2$ for mixed pieces.

\begin{lemma}[Section counts]\label{lem:qp:sections}
Fix a grid node $v$ of any level, a coordinate $i$ of $\gamma$, and real
values of the other coordinates. The set of values of $\gamma_i$ for which the
local event $E_v$ holds is a union of at most
$C_{\rm sec}=[(2k+1)P_{\rm sec}+1](8k+2)$ intervals.
\end{lemma}

\begin{proof}
Write $\gamma_i=t$. Then $\xi=\Psi\gamma$ and $\zeta=\Pi_T\gamma$ are affine in
$t$. Fix a query $a$. For every piece, membership of $a$ in its region is a
finite set of affine inequalities in $t$, hence an interval of $t$, and on
that interval the piece's value is a quadratic polynomial in $t$. For
continuous and separable pieces, every region certifies $W_\zeta(a)$ itself,
the regions cover the line (Lemmas~\ref{lem:qp:pieces} and~\ref{lem:qp:scalar}),
and valid formulas agree in value where intervals overlap. Hence between
consecutive interval endpoints, of which there are at most $2R$, a single
quadratic gives $W_\zeta(a)$. For mixed pieces, $W_\zeta(a)$ is the least value
of the pieces whose interval contains $t$: the pieces of one label give that
label's value, and the minimum over labels is $W_\zeta(a)$. Between
consecutive points among the at most $2R$ endpoints and the at most $R(R-1)$
crossings of distinct quadratics, the available pieces and their order are
constant, so again one quadratic gives the value; there are at most
$R^2+R\le2R^2$ such points. Gradients of overlapping formulas are not used.

The event $E_v$ involves at most $2k+1$ queries, giving at most
$(2k+1)P_{\rm sec}$ breakpoints. On each open interval between breakpoints, the
event is a conjunction of at most $2k$ inequalities between quadratics in $t$
(the term $\xi^{\mathsf T} a$ is affine in $t$). Each nonzero difference has at most
two roots, and identically zero differences have constant sign, so the
interval splits into at most $8k+1$ pieces on which the event is constant.
Counting the breakpoints themselves as possible isolated components gives at
most $[(2k+1)P_{\rm sec}+1](8k+1)+(2k+1)P_{\rm sec}\le C_{\rm sec}$ intervals.
\end{proof}

\begin{proof}[Proof of Lemma~\ref{lem:qp:volume}]
Let the columns of $V_0$ form a basis of $\ker T$ and $\mathsf S=[T^{\mathsf T}\ V_0]$,
which is invertible. Then $\gamma=\mathsf S(\xi,\eta)$ with $\zeta=V_0\eta$. The
vector $(\xi,\eta)$ has density $|\det\mathsf S|(2\sigma)^{-n}$ on
$\mathsf S^{-1}[-\sigma,\sigma]^n$. For fixed values of the coordinates other
than $\xi_Q$, the allowed set of $\xi_Q$ is contained in a box of volume
$\prod_{i\in Q}\ell_i$. The remaining coordinates range over the projection of
the parallelotope $\mathsf S^{-1}[-\sigma,\sigma]^n$, a zonotope whose volume is
$(2\sigma)^{n-t}\sum_{|J|=n-t}|\det(\mathsf S^{-1})_{Q^c,J}|$. By Jacobi's
complementary-minor identity,
$|\det\mathsf S|\,|\det(\mathsf S^{-1})_{Q^c,J}|=|\det\mathsf S_{J^c,Q}|=|\det T_{Q,J^c}|$.
Multiplying gives the first bound, with $K=J^c$. By Cauchy--Schwarz and the
Cauchy--Binet formula,
$\sum_{|K|=t}|\det T_{Q,K}|\le\binom nt^{1/2}\det(T_QT_Q^{\mathsf T})^{1/2}\le n^{t/2}$,
since $\norm{T_Q}\le\norm T\le1$.
\end{proof}

\begin{lemma}[Local counts under uniform ambient noise]\label{lem:qp:uniform-count}
Let $T$ have full row rank and $\norm T\le1$, let $\Lambda=\alpha I$, use
isotropic meshes on a box with widths $w_i$, and let $\gamma$ be continuous
uniform on $[-\sigma,\sigma]^n$. At every level,
$\sum_{v\in G_j}\Prob(E_v)\le\mathcal H_{\rm amb}=\prod_i[2+(1+2k)\alpha\sqrt nw_i/(2\sigma)]$.
\end{lemma}

\begin{proof}
For fixed $\zeta$, apply Lemma~\ref{lem:count:interval} to $W_\zeta$, which has
upper coordinate curvature $\alpha$ by Lemma~\ref{lem:qp:aux}(b). In the
comparisons of $E_v$ the term $\xi^{\mathsf T} a$ contributes only $\pm h_{ij}\xi_i$, so
$E_v$ confines each $\xi_i$, $i\in Q_v$, to an interval of length at most
$(1+2k)\alpha h_{ij}$ (Corollary~\ref{cor:count:levels}) that depends on
$\zeta$ only. Lemma~\ref{lem:qp:volume}, applied to the coordinates of $Q_v$
whose factor $\sqrt n(1+2k)\alpha h_{ij}/(2\sigma)$ is less than one, gives
$\Prob(E_v)\le\prod_{i\in Q_v}\min\{1,\sqrt n(1+2k)\alpha h_{ij}/(2\sigma)\}$.
Summing over the tensor grid as in the proof of Theorem~\ref{thm:count:local}
proves the claim.
\end{proof}

\begin{proof}[Proof of Lemma~\ref{lem:qp:gauss-count}]
\emph{Independence and a density majorant.} Under the proxy, $\xi=\Psi\gamma$
and $\zeta=\Pi_T\gamma$ are jointly Gaussian with
$\operatorname{Cov}(\xi,\zeta)=\sigma^2\Psi\Pi_T=\sigma^2(\Psi-\Psi T^{\mathsf T}\Psi)=0$,
because $\Psi T^{\mathsf T}=I$ and $\Pi_T$ is symmetric. They are therefore
independent, and
$\operatorname{Cov}\xi=\sigma^2(TT^{\mathsf T})^{-1}$ has eigenvalues in
$[\sigma^2,\sigma^2/c_{\rm fr}]$. Every principal submatrix of order $t$ has
eigenvalues in the same interval, so the density $p_Q$ of $\xi_Q$ satisfies
\[
 p_Q(y)\le(2\pi\sigma^2)^{-t/2}\exp\bigl(-c_{\rm fr}\norm y^2/(2\sigma^2)\bigr)
 =c_{\rm fr}^{-t/2}\prod_{i\in Q}\phi_s(y_i),
 \qquad s=\sigma/\sqrt{c_{\rm fr}},
\]
where $\phi_s$ is the centered normal density with standard deviation $s$.

\emph{Location.} Fix $v$ and $\zeta$, and let $x_v$ be an attaining witness of
$W_\zeta(v)$, chosen as a function of $\zeta$. With
$g=\Lambda(v-Tx_v)+\xi$, Lemma~\ref{lem:qp:aux}(c) gives
$V(v\pm h_{ij}e_i)\le V(v)\pm h_{ij}g_i+\lambda_ih_{ij}^2/2$. Together with the
comparisons of $E_v$ this yields $|g_i|\le\lambda_ih_{ij}/2+2E_j/h_{ij}\le C_k\lambda_ih_{ij}/2$
by Corollary~\ref{cor:count:levels}. Since $(Tx_v)_i\in[\ell_i,u_i]$, the event
$E_v$ implies $\xi_i\in J_i(v)=\lambda_i[\ell_i-v_i,u_i-v_i]+[-C_k\lambda_ih_{ij}/2,C_k\lambda_ih_{ij}/2]$,
a deterministic interval. As in Lemma~\ref{lem:qp:uniform-count}, $E_v$ also
confines $\xi_i$ to an interval $I_i(\zeta)$ of length at most
$C_k\lambda_ih_{ij}$. Only the event needs to lie in the intersection of the
two intervals.

\emph{Probability of one node.} By independence of $\xi$ and $\zeta$ and the
density majorant,
\[
 \Prob(E_v)\le c_{\rm fr}^{-|Q_v|/2}\prod_{i\in Q_v}
 \min\Bigl\{1,C_k\lambda_ih_{ij}\sup_{y\in J_i(v)}\phi_s(y)\Bigr\}.
\]
The factor $c_{\rm fr}^{-|Q_v|/2}$ stays outside the minima.

\emph{A weighted lattice sum.} Let $\Delta>0$, $C>0$, and $A\le B'$ with
$W'=B'-A$. We claim
\[
 \sum_{y\in y_0+\Delta\Z}\min\Bigl\{1,C\Delta\sup_{[A-y-C\Delta/2,\,B'-y+C\Delta/2]}\phi_s\Bigr\}
 \le CW'\phi_s(0)+2C+4 .
\]
At most $W'/\Delta+C+2$ lattice points satisfy
$y\in[A-C\Delta/2,B'+C\Delta/2]$, and each contributes at most
$\min\{1,C\Delta\phi_s(0)\}$, which gives at most $CW'\phi_s(0)+C+2$. For the
points beyond $B'+C\Delta/2$, the summand is
$\min\{1,C\Delta\phi_s(\theta)\}$, where $\theta$ is the distance to that
point. The summand decreases with $\theta$, so the tail sum is at most its
first term, bounded by $1$, plus its integral divided by $\Delta$.
This integral term is at most $C\int_0^\infty\phi_s=C/2$.
The tail contributes at most $1+C/2$; the same holds below $A-C\Delta/2$.

\emph{Summation.} Apply the claim with $\Delta=\lambda_ih_{ij}$, the lattice of
$\lambda_iv_i$ (the interior nodes of coordinate $i$ form a subset), core
$[\lambda_i\ell_i,\lambda_iu_i]$, and $C=C_k$. The two endpoint nodes
contribute one each. Expanding the product over the tensor grid,
\[
 \sum_{v\in G_j}\Prob(E_v)\le\prod_{i=1}^k\Bigl[2+c_{\rm fr}^{-1/2}\bigl(C_k\lambda_i(u_i-\ell_i)\phi_s(0)+2C_k+4\bigr)\Bigr],
\]
and $c_{\rm fr}^{-1/2}\phi_s(0)=1/(\sigma\sqrt{2\pi})$. Unrefined coordinates
have only their two endpoints. Nothing depends on the box.
\end{proof}

\subsection{Proofs of the main theorems}\label{app:qp:main}

We prove Theorems~\ref{thm:qp:gauss}, \ref{thm:qp:uniform},
and~\ref{thm:qp:aligned} together.

\emph{Base quantities.} If nonemptiness is not promised, first use
Proposition~\ref{prop:qp:primitives}(b), or (a) when $n_z=0$, with zero
objective to decide feasibility. If $X=\emptyset$, return that report
before constructing any thresholds, meshes, or law. On the remaining
nonempty instances, compute from the base data, in polynomial time, the
factorization (Lemma~\ref{lem:qp:normalize} or the supplied one, checked), the
ranges $[\ell_i,u_i]$, the bound $\Theta$ of Lemma~\ref{lem:qp:pieces}, the
integer bounds $L_i^z=\lceil\min_{\mathcal P}x_i\rceil$ and
$U_i^z=\lfloor\max_{\mathcal P}x_i\rfloor$ for $i\le n_z$, and
\[
 Z=\prod_i(U_i^z-L_i^z+1),\quad G=\sum_i(U_i^z-L_i^z),\quad R=Z2^r,\quad
 K=R(r+n+1),\quad B=\max\{2,Z2^r\},
\]
with $Z=1$ and $G=0$ if $n_z=0$; further $C_{\rm sec}$ from
Lemma~\ref{lem:qp:sections}, $\bar w$, $\theta_{\rm gap}$, and $C_{\rm gap}$ from
Lemma~\ref{lem:qp:gap}, and $\Psi=(TT^{\mathsf T})^{-1}T$ for the ambient laws.
All these numbers have bit length
polynomial in $I$, with absolute exponents; for example
$\log\Theta=O(n\log n+n\log C_0)$.

\emph{Sampling parameters.} Write $h_j=s2^{-j}$ for the nominal mesh on the box
$\mathcal A$ with largest width $s$, and $Q_J=(J+1)(2^J+1)^k$, which bounds the
number of nodes of all grids of levels $0,\dots,J$.
\begin{enumerate}[label=(\Alph*)]
\item \emph{Aligned law} ($n_z=0$, $\zeta=0$): $\mathcal A=\prod_i[\ell_i-\sigma/\alpha,u_i+\sigma/\alpha]$;
 $J$ is the least integer with $s2^{-J}\le\sigma/(2kKB(\alpha+\Theta))$; $M$
 is the least power of two with $M\ge\max\{2,2^J,2KB\}$.
\item \emph{Uniform ambient law}: $\mathcal A=\prod_i[\ell_i-s_i/\alpha,u_i+s_i/\alpha]$
 with $s_i=\sigma\norm{\Psi_{i,:}}_1$; $J$ is the least integer with
 $s2^{-J}\le\min\{1,\sigma/(2B[nkK(\alpha+\Theta)+GC_{\rm gap}])\}$; $M$ is the
 least power of two with $M\ge\max\{2,2B(K+G),2nC_{\rm sec}Q_J\}$.
\item \emph{Gaussian-like law}: for $t=0,1,2,\dots$ let
 $s_i(t)=2^t\sigma\norm{\Psi_{i,:}}_1$, let $\mathcal A_t$ and $s(t)$ be the
 corresponding box and largest width, let $J(t)$ be the least integer with
 $s(t)2^{-J}\le\min\{1,\sigma/(4B[kK(\alpha+\Theta)+GC_{\rm gap}])\}$, and let
 $b(t)$ be the least positive integer with
 $2^b\ge\max\{8B(nK+G),2nC_{\rm sec}Q_{J(t)}\}$. Stop at the first $t$ with
 $2^t\ge b(t)+20$, and use $\mathcal A=\mathcal A_t$, $J=J(t)$, $b=b(t)$.
\end{enumerate}
In (C), $s(t)\le2^ts(0)$, so $J(t)\le J(0)+t$, and $\log_2Q_{J(t)}\le
k(J(0)+t+1)+\log_2(J(0)+t+1)$. Since $J(0)$ and the other base quantities are
polynomial in $I$, $b(t)$ grows at most linearly in $t$ up to an additive
polynomial, and the loop stops after polynomially many steps with $J,b$
polynomial in $I$. In every case the sampled coefficients have bit length
polynomial in $I$, and the parameters are fixed before the draw.

\emph{Containment and correctness.} On every draw the box $\mathcal A$ contains
a global minimizer of $V$ by Lemma~\ref{lem:qp:aux}(d): in (A), $|\xi_i|\le\sigma$;
in (B), $|\xi_i|\le\sigma\norm{\Psi_{i,:}}_1$; in (C), the support of
$\mathcal G_{b,\sigma}$ lies in $[-\sigma(b+20),\sigma(b+20)]\subseteq[-2^t\sigma,2^t\sigma]$.
By Lemmas~\ref{lem:qp:pieces} and~\ref{lem:qp:gap}(b), every closure step of
the search of Definition~\ref{def:qp:closure-search} computes $\min_CV$ exactly
with an attaining witness. Theorem~\ref{thm:count:cells} therefore applies to
$V$ on $\mathcal A$; with Lemma~\ref{lem:qp:aux}(a) and
Lemma~\ref{lem:qp:face}(c), the output is exact on every draw, and its bit
length is polynomial.

\emph{Unresolved cells are rare.} Suppose a cell is retained and unclosed at
level $J$. By Theorem~\ref{thm:count:cells}(iv) it has a corner $v$ with
$V(v)-V^*\le2E_J$, and the closure test failed at $v$. Either some corner of the
cell lies outside $\mathcal R_{z_v,S_v}(\zeta)$, and then by
Lemma~\ref{lem:qp:tube} one of $K$ fixed functionals satisfies
$|\ell_\kappa|\le\eta_J=\sqrt k(\alpha+\Theta)h_J$; or $n_z\ge1$ and
$\Delta(v)<\theta_{\rm gap}h_J$, and then by Lemma~\ref{lem:qp:gap}(c), with
$h_J\le1$, the label gap of the sampled problem is at most $C_{\rm gap}h_J$.
This is one event about the sampled problem, not a union over visited cells.
By Lemmas~\ref{lem:qp:tube-prob} and~\ref{lem:qp:isolation}, the latter
applied conditionally on the continuous coefficients with $h(z)$ the slice
optimum, the probability that the fallback runs is at most
\[
 \begin{aligned}
 &\text{(A)}\quad K\bigl(k(\alpha+\Theta)h_J/\sigma+1/M\bigr)\le\tfrac1{2B}+\tfrac1{2B},\\
 &\text{(B)}\quad K\bigl(\sqrt{nk}(\alpha+\Theta)h_J/\sigma+1/M\bigr)+G\bigl(C_{\rm gap}h_J/\sigma+1/M\bigr)\le\tfrac1{2B}+\tfrac1{2B},\\
 &\text{(C)}\quad K\bigl(\eta_J/\sigma+2n2^{-b}\bigr)+G\bigl(C_{\rm gap}h_J/\sigma+2\cdot2^{-b}\bigr)\le\tfrac1{4B}+\tfrac1{4B}.
 \end{aligned}
\]
In (C) the concentration of $\mathcal G_{b,\sigma}$ is at most
$2\varepsilon/(\sigma\sqrt{2\pi})+2\cdot2^{-b}\le\varepsilon/\sigma+2\cdot2^{-b}$
on intervals of length $2\varepsilon$. The fallback costs
$B\poly(I)$, so by Lemma~\ref{lem:count:rare-fallback}(a) its expected
contribution is $\poly(I)$.

\emph{Expected local counts.} Let $N_j$ be the number of nodes of $G_j$ at which
the local event $E_v$ holds. By Theorem~\ref{thm:count:cells}(iv), at most
$2^kN_j$ cells are retained at level $j$, on every draw.
In (A), Corollary~\ref{cor:count:levels} with $M\ge2^J\ge m_{ij}$ gives
$\E N_j\le\mathcal H_{\rm al}$ for every $j\le J$.
In (B), Lemma~\ref{lem:qp:uniform-count} bounds the sum of the continuous
probabilities by $\mathcal H_{\rm amb}$ at every level. By
Lemma~\ref{lem:qp:sections} and Lemma~\ref{lem:count:transfer}(b) with
$\delta_i=1/M$, each node's probability changes by at most $2nC_{\rm sec}/M$,
and the choice of $M$ makes the total change over all levels at most one.
Hence $\sum_{j\le J}\E N_j\le(J+1)\mathcal H_{\rm amb}+1$.
In (C), Lemma~\ref{lem:qp:gauss-count} and the same transfer with
$\delta_i=2^{-b}$ give $\sum_{j\le J}\E N_j\le(J+1)\mathcal H_{\rm G}+1$.

\emph{Work.} Level $0$ processes one cell, and level $j+1$ processes at most
$2^k$ children of each cell retained at level $j$. A processed cell needs $2^k$
corner queries, each one exact solve (convex QP, or $2n_z+1$ convex MIQP solves
and a polishing solve) with polynomial bit length, a containment test of at
most $4^k(r+n)$ affine inequalities, and at most one closure by $3^k$ face
systems of order at most $k$. Its cost is at most
$f(n_z)\,6^k\poly(I)$. Summing, the expected work before the fallback is at
most $f(n_z)\,6^k\bigl(1+4^k\sum_{j\le J}\E N_j\bigr)\poly(I)$, and with
$J\le\poly(I)$ this gives the bounds of the three theorems.

\emph{Intrinsic bounds.} For the factorization of Lemma~\ref{lem:qp:normalize},
$u_i-\ell_i\le\norm T\diam(\mathcal P)\le\diam(\mathcal P)$ and $\alpha<4\nu$.
In Theorem~\ref{thm:qp:gauss}, $c_{\rm fr}^{-1/2}=(64/63)^{1/2}$ gives
$2+c_{\rm fr}^{-1/2}(4k+6)\le9+5k$, and $4/\sqrt{2\pi}\le2$. In
Theorem~\ref{thm:qp:uniform}, $\norm{(TT^{\mathsf T})^{-1}}\le64/63$ gives
$\norm{\Psi_{i,:}}_1\le\sqrt n\norm\Psi\le2\sqrt n$, hence
$w_i\le\diam(\mathcal P)+4\sigma\sqrt n/\alpha$. \qed

\subsection{The fiber estimate}

\begin{proof}[Proof of Proposition~\ref{prop:qp:fiber-sharp}]
Here $\xi=m\bar\gamma$, where $\bar\gamma$ is the mean of the coordinates, and
$\zeta=\gamma-\bar\gamma\mathbf1$. Given $\zeta$, the law of $\bar\gamma$ is
uniform on $[-\sigma-\min_l\zeta_l,\ \sigma-\max_l\zeta_l]$, and
$\xi-c(\zeta)=m(\min_l\gamma_l+\max_l\gamma_l)/2$ for the midpoint $c(\zeta)$
of the conditional support of $\xi$. Write $\gamma_l=\sigma(2U_l-1)$ with
independent uniform $U_l$ on $[0,1]$, and let $A',B'$ be their minimum and
maximum. The event is $|A'+B'-1|\le q$ with $q=\ell/(2\sigma m)$. Using the
joint density $n(n-1)(B'-A')^{n-2}$ for $n\ge2$ and integrating over the midpoint for a
fixed range $\varrho$ gives
$n(n-1)\int_0^1\varrho^{n-2}\min\{q,1-\varrho\}\,d\varrho=1-(1-q)^n$.
For $n=1$ we have $A'=B'=U_1$ and the interval
$\abs{2U_1-1}\le q$ has probability $q=1-(1-q)^1$, giving the same formula. As
$\ell\to0$ this is $nq(1+o(1))=\ell\sqrt n/(2\sigma)\,(1+o(1))$. Also
$\operatorname{Var}\xi=m^2\sigma^2/(3n)=\sigma^2/3$. For $k$ disjoint blocks the
events are independent, and each contributes the factor $m=(N/k)^{1/2}$, where
$N=km^2$ is the total dimension.
\end{proof}

Proposition~\ref{prop:qp:ambient-barrier} is the case $\alpha=1$ of
Theorem~\ref{thm:lim:ambient}; its proof is in Appendix~\ref{app:lim:ambient}.

\subsection{Separable recourse}\label{app:qp:sep}

\begin{proof}[Proof of Lemma~\ref{lem:qp:scalar}]
(a) Convexity of $\varphi_j$ makes $\delta_j$ nondecreasing. If
$\delta_j(t-1)\le\lambda\le\delta_j(t)$, then for $t'>t$,
$\varphi_j(t')-\varphi_j(t)=\sum_{s=t}^{t'-1}\delta_j(s)\ge(t'-t)\lambda$, and
symmetrically for $t'<t$; the converse follows by comparing $t$ with $t\pm1$.
Binary search for the least $t$ with $\delta_j(t)\ge\lambda$ (or the upper
endpoint if none exists) returns such a point. Each evaluation locates the
listed piece containing the integer and evaluates a rational quadratic.

(b) A point $t$ minimizes $\varphi_j(t)-\lambda t$ on the interval exactly
when $\lambda$ lies in the subdifferential of $\varphi_j$ at $t$, with the
normal cone included at the endpoints. In the interior of a piece with
$c_{jm}>0$ the subdifferential is $\{c_{jm}t+\mu_{jm}\}$; at a
breakpoint it is the interval between the one-sided derivatives. For every
$\lambda$ a minimizer exists. If it is a breakpoint or interior to a piece
with positive curvature, a listed state applies. If it is interior to a piece
with $c_{jm}=0$, then $\lambda=\mu_{jm}$ equals both one-sided derivatives at
the ends of that piece, and the breakpoint state at its left end applies.
\end{proof}

\begin{proof}[Proof of the separable piece properties and Theorem~\ref{thm:qp:sep}]
\emph{Pieces.} The inner objective is
$\frac\alpha2\norm a^2+\sum_j[\varphi_j(x_j)-\lambda_j(a,\zeta)x_j]$ with
$\lambda_j=\alpha t_j^{\mathsf T} a-\zeta_j$, so the scalar problems are independent,
and a choice $\omega$ of one state per coordinate whose intervals contain all
$\lambda_j$ certifies a global minimizer $x_\omega(a,\zeta)$, affine in
$(a,\zeta)$. Define $q_\omega$ by replacing each $\varphi_j$ by the listed
polynomial of its state (the constant $\varphi_j(t)$ for an integer state).
For a free coordinate, the first-order condition $\psi_{jm}'(x_j)=\lambda_j$
cancels the derivative of the response, and constant coordinates have zero
response derivative; hence
$\nabla_aq_\omega=\alpha(a-Tx_\omega)$ identically, and the Hessian is
$\alpha I-\alpha^2\sum_{\rm free}t_jt_j^{\mathsf T}/c_j$, bounded by
$\Theta_{\rm sep}$. The regions have at most $2n$ rows with normals
$\pm\alpha t_j^{\mathsf T}$ and offsets affine in $\zeta$, and they cover $\R^k$ for
every $\zeta$ by Lemma~\ref{lem:qp:scalar}. Every region certifies
$W_\zeta$ itself, so a cell inside a region is closed with no gap test, and
$P_{\rm sec}=2R_{\rm sep}$ in Lemma~\ref{lem:qp:sections}.

\emph{Fallback.} Each combination of an integer assignment and a listed piece
for every continuous coordinate gives a box on which $F_\gamma$ coincides with
a rational quadratic; minimize it by Lemma~\ref{lem:qp:face}(d) with $3^{n_c}$
faces and compare. The combinations cover $X$, so this is exact, at cost
$B_{\rm sep}\poly(I+L_{\rm in})$.

\emph{The theorem.} The proof of Appendix~\ref{app:qp:main} applies with the
following replacements: the exact recourse is the scalar procedure, of
polynomial cost and no factor $f(n_z)$; $R=R_{\rm sep}$, $K=R_{\rm sep}(2n+1)$,
$B=B_{\rm sep}$, $\Theta=\Theta_{\rm sep}$, $G=0$, and no gap test. In
Lemma~\ref{lem:qp:tube} the region rows are the scalar state inequalities, and
every step of its proof uses only the gradient identity, the attaining witness
at the query, and fixed normals with affine offsets. Part (a) uses law (A); integer
variables are allowed there because no gap event occurs. Part (b) uses law
(B).
\end{proof}

\begin{proof}[Proof of Lemma~\ref{lem:qp:aniso}]
Let $\Gamma=TT^{\mathsf T}\succ0$, let $H'\ge\norm\Gamma$ be its largest row sum, and
let $q_\Gamma$ clear the denominators of $\Gamma$. As in the proof of
Lemma~\ref{lem:qp:normalize}, $\mu'=q_\Gamma^{-k}H'^{-(k-1)}\le\lambda_{\min}(\Gamma)$.
Rational Jacobi rotations with entry threshold $\mu'/(64k)$ give an exactly
orthogonal rational $Q$ with $Q\Gamma Q^{\mathsf T}=\diag(a_i)+E$, $E$ with zero
diagonal and $\norm E\le\mu'/64$, in polynomial time. Each $a_i$ is a Rayleigh
quotient, so $\mu'\le a_i\le\norm T^2$. Choose dyadic $d_i$ with
$2a_i\le d_i^2<8a_i$, which exists because consecutive squared powers of two
differ by the factor four. With $\mathsf D=\diag(d_i)$ and
$U=\mathsf D^{-1}QT$, $UU^{\mathsf T}=\mathsf D^{-1}\diag(a_i)\mathsf D^{-1}+\mathsf D^{-1}E\mathsf D^{-1}$;
the first term has diagonal entries in $(1/8,1/2]$ and the second has norm at
most $(\mu'/64)/(2\mu')=1/128$, so $\frac1{16}I\preceq UU^{\mathsf T}\preceq I$. Finally
$U^{\mathsf T}\Lambda U=\alpha T^{\mathsf T} Q^{\mathsf T}\mathsf D^{-1}\mathsf D^2\mathsf D^{-1}QT=\alpha T^{\mathsf T} T$
because $Q$ is exactly orthogonal, and
$\lambda_i=\alpha d_i^2<8\alpha a_i\le8\beta$.
\end{proof}

\begin{proof}[Proof of Theorem~\ref{thm:qp:sep-gauss}]
By Lemma~\ref{lem:qp:aniso}, $F(x)=\sum_j\varphi_j(x_j)-\frac12(Ux)^{\mathsf T}\Lambda(Ux)$
exactly, so Lemma~\ref{lem:qp:aux} applies with $(U,\Lambda)$ in place of
$(T,\Lambda)$, and the scalar recourse applies with
$\lambda_j=(\Lambda u_j)^{\mathsf T} a-\zeta_j$, where $u_j$ is column $j$ of $U$. The
pieces satisfy $\nabla_aq_\omega=\Lambda(a-Ux_\omega)$, with Hessians
$\Lambda-\sum_{\rm free}(\Lambda u_j)(\Lambda u_j)^{\mathsf T}/c_j$ bounded by
$\Theta_\Lambda=\lambda_{\max}+\sum_{j\text{ continuous}}r_j^+\norm{\Lambda u_j}^2$.

Choose dyadic scales $\rho_i$ with $1\le\lambda_i\rho_i^2<4$ and use the
balanced nested meshes of Definition~\ref{def:count:mesh} with curvature
$\lambda$; then $c_{\rm bal}<4$, $C_k=1+8k$, and
$E_j\le\frac18\sum_i\lambda_i\rho_i^2h_j^2<kh_j^2/2$. In the proof of
Lemma~\ref{lem:qp:tube}, Lemma~\ref{lem:qp:aux}(c) now gives
$g^{\mathsf T}\Lambda^{-1}g\le4E_j<2kh_j^2$, so $\norm g\le\sqrt{2k\lambda_{\max}}\,h_j$,
and the cell diameter is at most $h_j\sum_i\rho_i$. The tube radius is
therefore at most $C_{\rm tube}h_j$ with the rational constant
$C_{\rm tube}=k(1+\lambda_{\max})+\Theta_\Lambda\sum_i\rho_i$. The ambient
normals have norm at least one because $\norm U\le1$. Lemma~\ref{lem:qp:gauss-count}
applies with $c_{\rm fr}=1/16$ and gives $\mathcal H_\Lambda$; its bound uses
$\lambda_i<8\beta$ and $u_i-\ell_i\le\norm U\diam(X)\le\diam(X)$.

For the law, use step (C) of Appendix~\ref{app:qp:main} with
$s_i(t)=2^t\sigma\norm{(A_U)_{i,:}}_1$, $A_U=(UU^{\mathsf T})^{-1}U$, the box
$\prod_i[\ell_i-s_i(t)/\lambda_i,u_i+s_i(t)/\lambda_i]$, the nominal scale
$S_\rho(t)=\max_i(u_i-\ell_i+2s_i(t)/\lambda_i)/\rho_i\le2^tS_\rho(0)$, the
cutoff $S_\rho(t)2^{-J}\le\sigma/(4KBC_{\rm tube})$, and
$2^b\ge\max\{8nKB,2nC_{\rm sec}Q_J\}$, where $K$, $B$, and $C_{\rm sec}$ are the
separable quantities and $G=0$. Then $J(t)\le J(0)+t$, the loop stops after
polynomially many steps, the probability of the fallback is at most
$K(C_{\rm tube}h_J/\sigma+2n2^{-b})\le1/(2B)$, and the rest of the proof is
unchanged.
\end{proof}
